%% kapa: exam=1 round=24 year=2013 period=1 form=A subject=economics subject_ko=경제원론 seq=14 number=54
%% source: Q-Net 감정평가사 기출문제 (cst00309) artlSeq=5073802 file=extracted/1008/5073802/2013년 제24회 감정평가사 1차시험 기출문제/1교시(A형).hwp
%% conversion: kordoc 4.24.51 -> tools/gen1.py (auto); review=proofread-sonnet; flags=none
\begin{question}{54}
독점기업이 8,000개의 상품을 판매하고 있다. 이때 상품가격은 1만 원, 평균총비용은 1만 2,000원, 평균가변비용은 8,000원이며, 한계비용과 한계수입은 6,000원으로 같다. 현재의 단기적인 상황에 관한 설명으로 옳은 것은?
\begin{choices}
  \item 총수입은 9,600만 원이다.
  \item 총비용은 8,000만 원이다.
  \item 상품 단위당 4,000원의 손실을 보고 있다.
  \item 생산을 중단하면 3,200만 원의 손실이 발생한다.
  \item 생산을 중단하는 것이 손실을 최소화한다.
\end{choices}
\end{question}
\answer{4}
